\documentclass[guides]{onlinebrief24}

% Absender (Zone 1)
\setreturnaddress{Max Mustermann, Musterstraße 1, 12345 Musterstadt}

% Empfänger (Zone 3)
\setrecipient{
  Erika Mustermann \\
  Musterweg 456 \\
  67890 Musterstadt
}

% Kopfdaten (werden nun manuell gesetzt, um Doppelungen zu vermeiden)
\setsubject{Testbrief mit exakter Fluchtlinie (25mm)}
\setplace{Musterstadt}
\setdate{\today}

\begin{document}

\begin{letter}{} % Leeres Argument, da wir \setrecipient nutzen

\opening{Sehr geehrte Frau Mustermann,}

dieser Brief demonstriert die exakte Einhaltung der Fluchtlinie bei 25mm. 

Sowohl die Absenderzeile (oben im Fenster), die Empfängeranschrift als auch dieser Text und der Betreff beginnen exakt bei 25mm vom linken Papierrand.

Das Adressfenster selbst (der Rahmen im `guides`-Modus) beginnt physisch bei 20mm, aber der Textinhalt ist um 5mm eingerückt.

Doppelte Datums- oder Betreffzeilen sollten nun der Vergangenheit angehören, da wir die automatische Ausgabe von KOMA-Script umgehen.

\closing{Mit freundlichen Grüßen}

\end{letter}

\end{document}
